\documentclass[11pt]{article}

\usepackage[margin=1in]{geometry}
\usepackage{amsmath,amssymb,amsthm}

\newtheorem{theorem}{Theorem}[section]
\newtheorem{lemma}[theorem]{Lemma}
\newtheorem{proposition}[theorem]{Proposition}
\newtheorem{corollary}[theorem]{Corollary}
\newtheorem{assumption}[theorem]{Assumption}
\theoremstyle{definition}
\newtheorem{definition}[theorem]{Definition}
\theoremstyle{remark}
\newtheorem{remark}[theorem]{Remark}

\newcommand{\area}{\operatorname{area}}
\newcommand{\dinv}{\operatorname{dinv}}
\newcommand{\defc}{\operatorname{defc}}
\newcommand{\UP}{\operatorname{UP}}
\newcommand{\DOWN}{\operatorname{DOWN}}
\newcommand{\East}{\operatorname{East}}
\newcommand{\West}{\operatorname{West}}
\newcommand{\inj}{\operatorname{inj}}

\title{The Length-Three and Length-Four Modulo-One Lower-Half Decomposition}
\author{}
\date{}

\begin{document}
\maketitle

\section{Statement and conventions}

Fix integers
\[
 \tau\ge2,\qquad 3\le s\le4,\qquad
 M_s=\tau\binom{s}{2},\qquad B_s=\tau(s-2)-1.
\]
A \emph{normalized step-$\tau$ Dyck word} is a word
\[
 D=(x_0,\ldots,x_{s-1})
\]
of nonnegative integers satisfying
\[
 x_0=0,\qquad x_{i+1}\le x_i+\tau.
\]
Its statistics are
\[
 \area(D)=\sum_i x_i,\qquad
 \dinv_\tau(D)=\sum_{i<j}d_\tau(x_i,x_j),\qquad
 \defc_\tau(D)=M_s-\area(D)-\dinv_\tau(D),
\]
where
\[
 d_\tau(u,v)=
 \begin{cases}
  (u+\tau-v)_+,&u\le v,\\
  (v+\tau+1-u)_+,&u>v.
 \end{cases}
\]
For a fixed defect $d$, put
\[
 L_d=\left\lfloor\frac{M_s-d}{2}\right\rfloor.
\]
Thus the weak lower half of the defect-$d$ layer is given by
$\area(D)\le L_d$.

An occurrence $x_j=e>0$ is \emph{extractable} if exactly one earlier
occurrence lies in
\[
 [\max(0,e-\tau),e)
\]
and deleting $x_j$ leaves a normalized word.  The leftmost extractable
occurrence is always chosen.  A normalized word with no extractable
occurrence is a \emph{full skeleton}.

Conversely, $\inj_e$ inserts $e$ immediately after the first occurrence in
$[\max(0,e-\tau),e)$ when the resulting word is normalized.  On the
configurations used below this reverses extraction.

On a three-window $(a,c,p)$, the map $\East_3$ is defined when
\[
 c\le p+\tau
\]
and is then the identity.  The map $\West_3$ is obtained by reversing the
window, applying $\East_3$, and reversing again; hence it is defined on
$(c,p,q)$ when
\[
 p\le c+\tau
\]
and is also the identity.  The one extraction preceding an East3 test may
occur in any nonfinal position.

\begin{assumption}[Root position]
\label{ass:root-position}
Every full skeleton $S$ of length $s$ with
$\defc_\tau(S)\le B_s$ satisfies
\[
 \area(S)\le L_{\defc_\tau(S)}.
\]
\end{assumption}

For $0\le d\le B_s$, let
\[
 \mathcal L_d
 =\{D:D\text{ is normalized},\ \defc_\tau(D)=d,\
                         \area(D)\le L_d\}
\]
and let $\mathcal F_d$ be the set of full skeletons of defect $d$.

\begin{theorem}[The $s=3,4$ lower-half decomposition]
\label{thm:main}
Under Assumption~\ref{ass:root-position}, there are mutually inverse maps
\[
 \UP:\{D\in\mathcal L_d:\area(D)<L_d\}\longrightarrow\mathcal L_d,
 \qquad
 \DOWN:\mathcal L_d\setminus\mathcal F_d\longrightarrow\mathcal L_d
\]
such that
\[
 \area(\UP(D))=\area(D)+1,\qquad
 \dinv_\tau(\UP(D))=\dinv_\tau(D)-1
\]
and
\[
 \area(\DOWN(D))=\area(D)-1,\qquad
 \dinv_\tau(\DOWN(D))=\dinv_\tau(D)+1.
\]
Consequently
\begin{equation}
 \mathcal L_d
 =\bigsqcup_{S\in\mathcal F_d}
   \{\UP^j(S):0\le j\le L_d-\area(S)\}.
 \label{eq:lower-union}
\end{equation}
In particular,
\begin{equation}
 \sum_{D\in\mathcal L_d}q^{\dinv_\tau(D)}t^{\area(D)}
 =
 \sum_{S\in\mathcal F_d}
 \sum_{j=\area(S)}^{L_d}q^{M_s-d-j}t^j.
 \label{eq:lower-gf}
\end{equation}
\end{theorem}

\section{Elementary extraction facts}

\begin{lemma}[Extraction transport]
\label{lem:transport}
Let $e$ be the leftmost extractable occurrence of $D$, let $C$ be obtained
by deleting it, and put $p=e-1$.  Then
\[
 \area(C:p)=\area(D)-1,\qquad
 \dinv_\tau(C:p)=\dinv_\tau(D)+1,\qquad
 \defc_\tau(C:p)=\defc_\tau(D).
\]
The identities hold even if the final adjacency of $C:p$ is invalid.
\end{lemma}

\begin{proof}
Write $D=A:e:E$.  For every $v\in E$,
\[
 d_\tau(e,v)=d_\tau(v,e-1).
\]
For $u\in A$, the difference
\[
 d_\tau(u,e-1)-d_\tau(u,e)
\]
is $1$ precisely when $u$ is in the predecessor interval of $e$, is $-1$
when $e<u\le e+\tau$, and is zero otherwise.  A leftmost extractable
occurrence is a record by the failed-splice propagation lemma in
\texttt{sgeq6\_decomp.tex}; that local proof is independent of the word
length.  Thus the negative case does not occur.  Extractability supplies
exactly one predecessor, so dinv increases by one.  Area decreases by one,
and defect is unchanged.
\end{proof}

\begin{lemma}[The small full skeletons]
\label{lem:full-classification}
The full skeletons have the following forms:
\[
 \begin{array}{c|c}
 s=3&(0,0,a),\quad 0\le a\le\tau,\\[1mm]
 s=4&(0,0,a,b),\quad 0\le a,b\le\tau.
 \end{array}
\]
For a full skeleton of defect at most $B_s$, the full-skeleton UP operation
\begin{equation}
 S=C:a\longmapsto\inj_{a+1}(C)
 \label{eq:full-up}
\end{equation}
is defined and normalized.
\end{lemma}

\begin{proof}
If a normalized word begins $(0,u)$ with $u>0$, then $u$ has the initial
zero as its unique predecessor.  If its deletion splice fails, the next
entry is a new record with $u$ as unique predecessor.  Repeating reaches an
extractable final occurrence.  Thus every full word begins $(0,0)$.

At length three, the last entry is at most $\tau$ and hence has both zeros
as predecessors.  This proves the first line.  At length four, write the
word as $(0,0,a,b)$.  We have $a\le\tau$.  If $b>\tau$, normalization makes
$a$ its unique predecessor, so $b$ is extractable at the right boundary.
Thus $b\le\tau$.  The converses follow because every positive entry then
has both initial zeros as predecessors.

For $s=3$, the only full word for which $a+1$ has no anchor in $C=(0,0)$
is $(0,0,\tau)$.  Its defect is $\tau=B_3+1$.  For $s=4$, an anchor can
fail only for $(0,0,0,\tau)$; if the penultimate entry is positive, it
anchors $\tau+1$.  The exceptional word has defect $2\tau=B_4+1$.
Every other case in the stated defect range therefore has the required
anchor, and the adjacent inequalities give normalization.
\end{proof}

\begin{lemma}[The small full-prefix DOWN branch]
\label{lem:full-prefix}
Suppose the leftmost extraction of a word $D$ of defect at most $B_s$ has
value $e=p+1$ and leaves a full word $C$.  Then $C:p$ is normalized and
full.
\end{lemma}

\begin{proof}
For $s=3$, necessarily $C=(0,0)$.  The extracted value must have occurred
between the zeros, so $p\le\tau-1$, and $(0,0,p)$ is full.

For $s=4$, write $C=(0,0,c)$ with $c\le\tau$.  If the extraction was
nonfinal, its position and normalization give $p\le\tau-1$.  If it was
final, then $e\le c+\tau$, so $C:p$ is normalized.  Were $p>\tau$, the
word $(0,0,c,p)$ would have
\[
 \dinv_\tau(0,0,c,p)=4\tau-c-p,\qquad
 \defc_\tau(0,0,c,p)=2\tau.
\]
This contradicts transport and $\defc_\tau(D)\le B_4=2\tau-1$.  Hence
$p\le\tau$, and the two initial zeros show that $C:p$ is full.
\end{proof}

\begin{lemma}[A nonfinal first UP extraction]
\label{lem:up-nonfinal}
If $D$ has defect at most $B_s$, is strictly below the middle of its defect
layer, and is not full, then its leftmost extraction is nonfinal.
\end{lemma}

\begin{proof}
Strictly below the middle implies
\begin{equation}
 \dinv_\tau(D)-\area(D)\ge2.
 \label{eq:strict-gap}
\end{equation}
Suppose the only extractable occurrence is final.

For $s=3$, a word beginning $(0,0)$ cannot have an extractable final
occurrence.  Otherwise write $D=(0,a,c)$ with $a>0$.  The failure of the
deletion splice for $a$ gives
\[
 1\le a\le\tau,\qquad \tau<c\le a+\tau.
\]
Only adjacent pairs contribute, and
\[
 \dinv_\tau(D)-\area(D)=2\tau-a-2c<0,
\]
contrary to~\eqref{eq:strict-gap}.

For $s=4$, first suppose $D=(0,0,a,c)$.  A final extraction forces
$0<a\le\tau<c\le a+\tau$, and direct substitution gives
\[
 \defc_\tau(D)=2\tau>B_4.
\]
If instead $D=(0,a,b,c)$ with $a>0$, the first two nonfinal occurrences
are not extractable only if
\[
 b>\tau,\qquad c>a+\tau.
\]
The word is then a rapid record chain, only adjacent pairs contribute, and
\[
 \dinv_\tau(D)=3\tau-c<\area(D).
\]
This again contradicts~\eqref{eq:strict-gap}.
\end{proof}

\begin{lemma}[Ordinary level-three reconstruction]
\label{lem:level-three}
Suppose the first UP extraction is nonfinal and leaves $C=A:c$.  Put
$p=e-1$.  If $c\le p+\tau$, then
\begin{equation}
 \inj_{c+1}\bigl(\inj_{p+1}(A)\bigr)
 \label{eq:east3-up}
\end{equation}
is defined and normalized.  Conversely, after two successive DOWN
extractions, a successful West3 stage is reconstructed by
\begin{equation}
 C:p:q\longmapsto\inj_{q+1}(C:p),
 \label{eq:west3-down}
\end{equation}
and the two operations are inverse.
\end{lemma}

\begin{proof}
Because the first extraction was nonfinal, injecting $p+1$ into $A$
recovers the original word without its last occurrence $c$.  If every
entry of this prefix is at most $c$, the inequality $c\le p+\tau$ makes
$p+1$ an anchor for $c+1$.  If some entry exceeds $c$, the predecessor of
the first such entry lies in $[c+1-\tau,c]$ and is an anchor.  In either
case insertion after the first anchor is normalized.

For the reverse direction, write the two extracted values as $p+1,q+1$.
The West3 condition $p\le C[-1]+\tau$, together with the splice and anchor
inequalities belonging to the two extractions, makes~\eqref{eq:west3-down}
normalized.  The inserted occurrence is extracted first, after which the
East3 condition is exactly the displayed West3 inequality.  This reverses
\eqref{eq:east3-up}.  This is the same level-three check used in
\texttt{sgeq6\_decomp.tex} and \texttt{s5\_decomp.tex}; importantly, it
requires only that the UP extraction be nonfinal.
\end{proof}

\section{The length-three tests}

\begin{lemma}[East3 and West3 do not fail for $s=3$]
\label{lem:s3-total}
At length three, every East3 or West3 test arising in the lower-half
algorithm is successful.
\end{lemma}

\begin{proof}
For UP, Lemma~\ref{lem:up-nonfinal} makes the first extraction the middle
entry.  Its deletion leaves $(0,c)$ normalized, so $c\le\tau\).  The East3
stage is $(0,c,p)$, and
\[
 c\le\tau\le p+\tau.
\]

For DOWN, suppose the first extraction leaves a nonfull word, so a second
extraction is required.  If the first extraction had been final, it would
have been the only extraction of the original word.  The calculation in
the $s=3$ part of Lemma~\ref{lem:up-nonfinal} gives
$\dinv_\tau(D)<\area(D)$, contradicting even the weak lower-half condition.
Thus the first extraction is the middle entry and its lowered value
$p$ satisfies $p\le\tau-1$.  After the second extraction the West3 stage
is $(0,p,q)$, so its condition $p\le\tau$ holds.
\end{proof}

\section{The length-four boundary pair}

\begin{definition}[The boundary pair]
\label{def:boundary}
For integers satisfying
\begin{equation}
 0\le p<a\le\tau-1,\qquad p+\tau<c\le a+\tau,
 \label{eq:boundary-bounds}
\end{equation}
define
\begin{equation}
 \mathcal B_\uparrow(0,p+1,a,c)=(0,a+1,c+1,p),
 \qquad
 \mathcal B_\downarrow(0,a+1,c+1,p)=(0,p+1,a,c).
 \label{eq:boundary-pair}
\end{equation}
\end{definition}

\begin{lemma}[Classification of failed East3]
\label{lem:east-failure}
Let $s=4$, and suppose $D$ satisfies the UP hypotheses.  If its East3 test
fails, then
\[
 D=(0,p+1,a,c)
\]
with the bounds~\eqref{eq:boundary-bounds}.
\end{lemma}

\begin{proof}
By Lemma~\ref{lem:up-nonfinal}, the extracted value $p+1$ is nonfinal.
Write the remainder as $(0,a,c)$.  The extracted occurrence cannot have
been in position two.  Indeed, East3 failure and the original final
adjacency would then give $c=(p+1)+\tau$, while the deletion splice would
give $p+1\le a$.  The earlier occurrence $a$ would consequently be
extractable before $p+1$, a contradiction.  Hence
\[
 D=(0,p+1,a,c).
\]
The deletion splice gives $a\le\tau$, normalization gives
$c\le a+\tau$, and East3 failure gives $c>p+\tau$.  Therefore $p<a$ and
$p\le\tau-1$.

It remains to exclude $a=\tau$.  In that case direct substitution gives
\[
 \dinv_\tau(D)-\area(D)=2\tau-p-1-2c<0,
\]
because $c>p+\tau$.  This contradicts strict lower-half position.  Thus
$a\le\tau-1$, proving all the bounds.
\end{proof}

\begin{lemma}[Classification of failed West3]
\label{lem:west-failure}
Let $s=4$, and suppose $D$ is weakly below the middle and its DOWN
algorithm reaches a failed West3 stage.  Then
\[
 D=(0,a+1,c+1,p)
\]
with the bounds~\eqref{eq:boundary-bounds}.
\end{lemma}

\begin{proof}
Let the two successive extracted values be $c+1$ and $a+1$, and write the
twice-deleted word as $(0,p)$.  Thus the West3 stage is
\[
 (0,p,c,a),
\]
and failure says $c>p+\tau$.

Neither $0$ nor $p$ can be the unique predecessor of $c+1$.  Hence the
second extracted value $a+1$ occurs before $c+1$, is its unique predecessor,
and immediately precedes it.  In particular
\[
 c\le a+\tau.
\]
There are two possible orders for the remaining occurrence $p$:
\[
 (0,p,a+1,c+1)
 \quad\hbox{or}\quad
 (0,a+1,c+1,p).
\]
In the first order the only extraction in the original word is the final
entry $c+1$: deletion of $a+1$ is blocked by
$c+1>p+\tau$, and all earlier entries precede the chosen leftmost
extraction.  The rapid-chain calculation from
Lemma~\ref{lem:up-nonfinal} gives
$\dinv_\tau(D)<\area(D)$, contradicting weak lower-half position.
Thus the second order holds.

Now $a+1$ occurs in position one of its once-deleted word, so
$a\le\tau-1$.  Finally, $c\le a+\tau$ and $c>p+\tau$ imply $p<a$.
These are exactly~\eqref{eq:boundary-bounds}.
\end{proof}

\begin{lemma}[Boundary normalization, statistics, and recognition]
\label{lem:boundary-check}
The two words in~\eqref{eq:boundary-pair} are normalized.  The boundary
maps are inverse, and
\[
 (\area,\dinv_\tau)\bigl(\mathcal B_\uparrow(D)\bigr)
 =(\area(D)+1,\dinv_\tau(D)-1).
\]
Moreover, DOWN recognizes the image of $\mathcal B_\uparrow$ by extracting
$c+1$ and then $a+1$, after which West3 fails; UP recognizes the image of
$\mathcal B_\downarrow$ by extracting $p+1$, after which East3 fails.
\end{lemma}

\begin{proof}
For the UP image, normalization follows from
\[
 a+1\le\tau,\qquad c+1\le a+1+\tau.
\]
The inverse image is normalized because $p+1\le\tau$,
$a\le p+1+\tau$, and $c\le a+\tau$.

The zero-pair contributions agree on the two sides:
\begin{align*}
 &d_\tau(0,p+1)+d_\tau(0,a)+d_\tau(0,c)\\
 &\qquad =
 d_\tau(0,a+1)+d_\tau(0,c+1)+d_\tau(0,p).
\end{align*}
Also
\[
 d_\tau(a,c)=d_\tau(a+1,c+1),\qquad
 d_\tau(p+1,c)=d_\tau(c+1,p)=0.
\]
The remaining pair loses exactly one:
\[
 d_\tau(p+1,a)=d_\tau(a+1,p)+1.
\]
The area visibly increases by one, proving the statistic assertion.

In $(0,a+1,c+1,p)$, the entry $a+1$ is not extractable initially because
deleting it exposes $c+1>\tau$.  The entry $c+1$ has $a+1$ as its unique
predecessor and has a valid deletion splice, so it is extracted first.
Then $a+1$ is extracted, leaving $(0,p)$ and producing the failed West3
stage $(0,p,c,a)$.  Conversely, in $(0,p+1,a,c)$ the occurrence $p+1$ is
extractable, and deletion gives $(0,a,c)$ and the failed East3 stage
$(0,a,c,p)$.  This proves branch recognition and invertibility.
\end{proof}

\section{Definition and verification of the maps}

Let $D$ have defect at most $B_s$ and be strictly below the middle.  Define
$\UP$ as follows.
\begin{enumerate}
\item If $D=C:a$ is full, use the operation~\eqref{eq:full-up}.
\item Otherwise extract $e=p+1$, leaving $C=A:c$.  If East3 succeeds, use
the ordinary operation~\eqref{eq:east3-up}.
\item If East3 fails, then $s=4$ by Lemma~\ref{lem:s3-total}; use
$\mathcal B_\uparrow$ from Definition~\ref{def:boundary}.
\end{enumerate}

Let $D$ have defect at most $B_s$, lie weakly below the middle, and not be
full.  Define $\DOWN$ as follows.
\begin{enumerate}
\item Extract $e_1=p+1$, leaving $C_1$.  If $C_1$ is full, return $C_1:p$.
\item Otherwise extract $e_2=q+1$, leaving $C_2$, and form $C_2:p:q$.
If West3 succeeds, use the ordinary operation~\eqref{eq:west3-down}.
\item If West3 fails, then $s=4$ by Lemma~\ref{lem:s3-total}; use
$\mathcal B_\downarrow$ from Definition~\ref{def:boundary}.
\end{enumerate}

\begin{proposition}[Totality, statistics, and invertibility]
\label{prop:maps}
The maps are normalized and total on the stated domains.  They preserve
defect, change $(\area,\dinv_\tau)$ by $(+1,-1)$ for UP and by $(-1,+1)$
for DOWN, and are inverse wherever both sides are defined.
\end{proposition}

\begin{proof}
Lemmas~\ref{lem:full-classification} and~\ref{lem:up-nonfinal} give the
first UP extraction and the full branch.  Lemma~\ref{lem:s3-total} handles
all remaining length-three tests.  At length four, an East3 failure is
exactly the boundary case by Lemma~\ref{lem:east-failure}.

For DOWN, Lemma~\ref{lem:full-prefix} gives the full-prefix branch.  If the
remainder is not full, a second extraction exists by definition.
Lemma~\ref{lem:s3-total} handles length three, while
Lemma~\ref{lem:west-failure} classifies every failed length-four West3
stage.  Ordinary branch normalization is Lemma~\ref{lem:level-three}, and
boundary normalization is Lemma~\ref{lem:boundary-check}.  Thus the maps
are total.

In an ordinary UP branch, transport changes statistics by $(-1,+1)$ and
the two raised inverse injections change them by $(+2,-2)$, for a net
change $(+1,-1)$.  The ordinary DOWN branch has two transports and one
inverse injection, for a net change $(-1,+1)$.  The full branches are the
same counts with zero and one extraction, respectively.  The boundary
statistics were computed in Lemma~\ref{lem:boundary-check}.

The full branches reverse one another because the inserted occurrence is
the first extraction and leaves the same full prefix.  The ordinary
East3 and West3 branches reverse one another by
Lemma~\ref{lem:level-three}.  The two boundary branches reverse and
recognize one another by Lemma~\ref{lem:boundary-check}.  Hence the maps
are inverse.  Defect preservation follows from
$M_s-\area-\dinv_\tau$.
\end{proof}

\section{Proof of the lower-half decomposition}

\begin{proof}[Proof of Theorem~\ref{thm:main}]
Fix $0\le d\le B_s$.  Proposition~\ref{prop:maps} defines UP at every area
below $L_d$ and DOWN at every weakly lower-half nonfull word.  Starting at
any word of $\mathcal L_d$ and iterating DOWN must terminate because area
decreases at every step.  It can terminate only at a full skeleton.
Invertibility makes the root and the entire descent unique.

Conversely, Assumption~\ref{ass:root-position} places each
$S\in\mathcal F_d$ in $\mathcal L_d$.  Iterating UP gives exactly one word
at each area
\[
 \area(S),\area(S)+1,\ldots,L_d.
\]
No two strings meet, again by invertibility.  This proves
\eqref{eq:lower-union}.  A word of defect $d$ and area $j$ has dinv
$M_s-d-j$, so summing the strings proves~\eqref{eq:lower-gf}.
\end{proof}

\end{document}
